\begin{titlepage}
    \centering
    \vspace*{2.6cm}

    {\zihao{1}\bfseries \ReportTitle \par}
    \vspace{0.9cm}
    {\zihao{3}\bfseries 面向七星微赛题的工程化 RISC-V 处理器方案 \par}
    \vspace{1.6cm}

    {\zihao{2}\bfseries 参赛作品：\ProjectName \par}
    \vspace{1.4cm}

    \renewcommand{\arraystretch}{1.35}
    \begin{tabular}{rl}
        项目名称： & \ProjectName \\
        报告版本： & \ReportVersion \\
        更新日期： & \ReportDate \\
        冻结主口径： & \texttt{RV32I + Zicsr} \\
        目标原型平台： & Nexys A7-100T \\
        文档性质： & 初赛正式设计说明书 \\
    \end{tabular}

    \vfill
\end{titlepage}

\section*{摘要}
\addcontentsline{toc}{section}{摘要}

\texttt{YH\_rv\_cpu} 是一套面向七星微赛题的工程化 RISC-V 处理器方案。该设计以
\texttt{RV32I + Zicsr} 为冻结性能与提交主口径，采用 IF/ID/EX/MEM/WB 五级顺序流水，
并构建了由同步指令存储、数据存储、UART、timer 和 DONE 标志组成的最小 SoC，
使 RTL、软件、回归验证、CoreMark 和 FPGA 原型路径能够保持一致。

当前工程已经形成覆盖功能正确性、性能评估和原型实现可行性的完整证据链。功能验证方面，
除基础 smoke 回归外，还完成了 \texttt{rv32 baseline 33/33}、
\texttt{rv64 baseline 21/21}、\texttt{rv32 full-ui 42/42} 和
\texttt{rv64 full-ui 54/54}。性能方面，CoreMark short 为
\texttt{0.912472 CoreMark/MHz}，strict 长跑为
\texttt{0.912465 CoreMark/MHz} 且运行时间 \texttt{10.959325s}。
原型实现方面，Vivado \texttt{impl50} 结果为
\texttt{2556 LUT / 2170 FF / 4 BRAM / 0 DSP}，并保持
\texttt{WNS = +5.599ns}、\texttt{WHS = +0.025ns} 的正时序裕量。

综合来看，该设计已经具备清晰的微架构边界、可复核的验证矩阵、稳定的性能记录以及明确的 FPGA 原型适配路径。
同时，当前工程还具备 \texttt{RV32/RV64} 双 XLEN 验证能力，为后续实体板 bring-up、
控制流优化和系统扩展保留了清晰起点。

\section*{关键词}
\addcontentsline{toc}{section}{关键词}

RISC-V；五级流水；工程化验证；CoreMark；FPGA 原型；双 XLEN

\section*{重要内容预览}
\addcontentsline{toc}{section}{重要内容预览}

\begin{enumerate}
    \item \textbf{赛题要求对照总表}：正文先给出“赛题要求-本文亮点-当前状态”对应关系，不把题面中的条目当成可忽略项，而是逐条说明当前已覆盖、部分覆盖和待补齐内容。
    \item \textbf{CPU 架构设计说明书}：正文给出五级流水 CPU、最小 SoC、冒险处理、CSR/trap/timer 支持与模块职责划分，并明确哪些 IF 前端增强项已落地、哪些仍是后续对齐赛题的方向。
    \item \textbf{Verilog/VHDL 建模规范}：正文单列 HDL 建模规范章节，说明模块划分、组合/时序逻辑约束、参数化策略、接口组织方式与关键模块交互信号组。
    \item \textbf{验证计划与测试报告}：正文给出 smoke、baseline、扩展 \texttt{full-ui}、CoreMark、Vivado \texttt{impl50} 与 FPGA-like probe 的验证矩阵和主要结果，并明确简单应用与演示材料的当前边界。
    \item \textbf{FPGA 适配与后续补齐项}：正文说明目标板卡、固定原型参数、综合实现口径、报告位置以及后续实体板、交付包和演示材料的补齐路径。
\end{enumerate}

\begin{table}[H]
    \centering
    \small
    \caption{赛题要求与本文亮点总览}
    \begin{tabularx}{0.95\textwidth}{L{0.25\textwidth} Y L{0.15\textwidth}}
        \toprule
        赛题关注点 & 本文亮点与当前口径 & 当前状态 \\
        \midrule
        RV32I/RV64I CPU 与扩展边界 & 以 \texttt{RV32I + Zicsr} 为冻结提交口径，同时保留 \texttt{RV32/RV64} 双 XLEN 验证结果与后续扩展边界说明 & 已覆盖 \\
        五级流水与三类冒险 & 正文给出 IF/ID/EX/MEM/WB 组织、前递、load-use 最小停顿、redirect 与 trap 统一控制流框架 & 已覆盖 \\
        HDL 建模与模块化实现 & 正文单列建模规范，说明模块职责、参数化策略和关键模块交互信号组 & 已覆盖 \\
        分层仿真、性能测试与报告 & 正文汇总 smoke、\texttt{riscv-tests}、CoreMark、\texttt{impl50} 与 FPGA-like probe 结果 & 已覆盖 \\
        FPGA 原型与 50MHz 路径 & 正文给出 Nexys A7-100T、Vivado \texttt{impl50}、资源和时序结果，并说明 pre-board closure 边界 & 部分覆盖 \\
        简单应用、演示视频与完整交付包 & 正文明确这部分仍需继续补齐，避免把已完成验证误写成赛题全部闭环 & 待补齐 \\
        \bottomrule
    \end{tabularx}
\end{table}

\clearpage
